\section{Finite-index discrepancy estimates from the actual asymptotic order}

\begin{lemma}[An original member attains the actual mapped DVR ideal order]
\label{lem:c0n-original-member-attains-mapped-dvr-order}
\lean{EnvelopeProperty.Char0.exists_mem_eq_dvrIdealOrder_map}
\leanok
Let $R$ be an original commutative ring, $V$ a discrete valuation ring,
and $g:R\to V$ the original ring homomorphism.  If the actual extension
$g_*I$ of an original ideal $I\subset R$ is nonzero, then an original
$x\in I$ satisfies $x\ne0$ and
$\operatorname{ord}_V(g(x))=\operatorname{ord}_V(g_*I)$.
No injectivity of $g$ is assumed.
\end{lemma}
\begin{proof}
This reuses the proved native original-member theorem.  Write
$r=\operatorname{ord}_V(g_*I)$, which is a finite natural number.
If every original member had image of order greater than $r$, all those
images would lie in the $(r+1)$st maximal-ideal power.  Their whole
extended ideal would lie there too.  The actual DVR order reverses ideal
containment, and that power has order $r+1$, giving $r+1\le r$.
Some original member therefore has image of order at most $r$.
Membership in the same extended ideal gives the opposite inequality,
so its order is exactly $r$.  Finite image order makes its image nonzero,
and consequently the original member itself is nonzero.
\end{proof}

\begin{lemma}[A finite index approximates the normalized-order infimum]
\label{lem:c0n-finite-index-infimum}
\uses{def:sc0-original-asymptotic-divisorial-order}
\lean{EnvelopeProperty.Char0.exists_positive_index_normalized_order_lt}
\leanok
For any nonnegative real sequence $N_q$ and any $m\in\mathbf N$, put
$\alpha=\inf_{q>0}N_q/q$.  There is $q>0$ with
\[
 (m/q)N_q<m\alpha+1.
\]
\end{lemma}
\begin{proof}
For $m=0$ take $q=1$.  For $m>0$, the definition of infimum provides
$q>0$ with $N_q/q<\alpha+1/m$.  Multiplication by $m$ gives the assertion.
\end{proof}

\begin{theorem}[Actual finite-index element discrepancy bound]
\label{thm:c0n-actual-finite-index-discrepancy}
\uses{lem:c0n-finite-index-infimum,thm:sc0-original-asymptotic-multiplier-divisorial-bound}
\lean{EnvelopeProperty.Char0.exists_asymptoticMultiplierIdeal_element_order_bound}
\leanok
Let $k$ be a field of characteristic zero, let $K/k$ be a field extension,
and let $C,V\subset K$ be $k$-subalgebras with fraction field $K$.
Assume that $C$ is Noetherian, that $V$ is an essentially finite type
discrete valuation ring over $k$, and that $C\subset V$.
Let $J_\bullet$ be a graded sequence in $C$ with $J_q\ne0$ for $q>0$.
Write $B_m=\mathcal J(J_\bullet^m)$ for the full asymptotic multiplier ideal,
$N_q=\operatorname{ord}_V(J_q\mathcal O_V)$, and
$j=\operatorname{ord}_V(\operatorname{Jac}_{C,V})$ for the intrinsic
original-source Jacobian order.  For every $m\in\mathbf N$ there is $q>0$
such that every $f\in B_m$ of finite order $\nu$ satisfies
\[
 (m/q)N_q-(2+j)\le\nu.
\]
The index $q$ is independent of the chosen element $f$.
\end{theorem}
\begin{proof}
Choose $q$ from the preceding lemma for the actual normalized-order
infimum.  The original asymptotic multiplier discrepancy estimate gives
$m\alpha<\nu+1+j$ for every nonzero $f\in B_m$.  An element of finite
order is nonzero.  Combining the two strict inequalities gives
$(m/q)N_q<\nu+2+j$, and hence the displayed weak inequality.
This argument also applies when $m=0$.
\end{proof}

\begin{corollary}[All lower orders of the original sequence ideal]
\label{cor:c0n-actual-lower-orders}
\uses{thm:c0n-actual-finite-index-discrepancy,lem:c0n-original-member-attains-mapped-dvr-order}
\lean{EnvelopeProperty.Char0.exists_asymptoticMultiplierIdeal_lower_order_bound}
\leanok
For the same $q$, every $o\in\mathbf N$ with
$o\le\operatorname{ord}_V(g)$ for all $g\in J_q$ satisfies
$(m/q)o-(2+j)\le\operatorname{ord}_V(f)$ for every finite-order
$f\in B_m$.
\end{corollary}
\begin{proof}
The finite order $N_q$ of the extended nonzero ideal is attained by the
image of an original element of $J_q$.  Thus $o\le N_q$.
Since $m/q\ge0$, the asserted bound follows from the preceding theorem.
\end{proof}

\begin{corollary}[Extended-real logarithmic bound, including zero]
\label{cor:c0n-actual-ereal-log-bound}
\uses{thm:c0n-actual-finite-index-discrepancy}
\lean{EnvelopeProperty.Char0.exists_asymptoticMultiplierIdeal_ereal_order_bound}
\leanok
Fix a real number $a>0$.  For every $m\in\mathbf N$ there is $q>0$ such
that all $f\in B_m$, including $f=0$, satisfy
\[
 -a\operatorname{ord}_V(f)\le
 (m/q)(-aN_q)+a(2+j)
\]
in the extended real numbers.
\end{corollary}
\begin{proof}
Use the index of the finite-order theorem.  For $f\ne0$, multiply its
finite-order estimate by the negative number $-a$ and rearrange.
For $f=0$, the left side is $-\infty$, so the inequality holds.
The finite order $N_q$ is legitimate because the original ideal $J_q$
is nonzero and its inclusion into the divisorial ring is injective.
\end{proof}

\begin{corollary}[The same estimate for original affine sections]
\label{cor:c0n-original-scheme-ereal-log-bound}
\uses{cor:c0n-actual-ereal-log-bound,thm:sc0-original-global-asymptotic-affine-values}
\lean{EnvelopeProperty.Char0.exists_schemeAsymptoticMultiplierIdealSheaf_ereal_order_bound}
\leanok
Let $Y$ be a smooth integral scheme over $k$, let $J_\bullet$ be a graded
ideal-sheaf sequence whose positive-index generic stalks are nonzero,
and let $U\subset Y$ be an original nonempty affine open.  For every
divisorial discrete valuation ring $V$ of the actual function field of $Y$
containing the image of $\Gamma(U,\mathcal O_Y)$, the preceding
extended-real estimate holds for every $m\in\mathbf N$ and $a>0$ and all
sections of the full original
asymptotic multiplier ideal on $U$.  The order of $J_q$ is computed using
the actual composite from original affine sections into $V$.
\end{corollary}
\begin{proof}
Identify the original affine section ring with its image in the actual
function field.  Smoothness supplies Noetherianity of this source algebra,
and the nonzero generic stalk hypothesis supplies nonzero positive-index
affine ideals.  The exact original affine comparison identifies the
whole scheme asymptotic ideal with the inverse image of the full affine
asymptotic ideal under this ring isomorphism.  Apply the preceding
extended-real estimate and compose the original section map with the
inclusion into $V$.  Extension of ideals under this composite equals
successive extension under the two maps, so both the section orders and
the original sequence ideal orders are exactly those asserted.
\end{proof}
